\documentclass[11pt]{article}
\usepackage[margin=1in]{geometry}
\usepackage[T1]{fontenc}
\usepackage[utf8]{inputenc}
\usepackage{lmodern}
\usepackage{amsmath,amssymb,amsthm,mathtools}
\usepackage{booktabs,tabularx,longtable,array}
\usepackage{xltabular}% tabularx-style tables that can break across pages
\usepackage{enumitem}
\usepackage[dvipsnames]{xcolor}
\usepackage{cite}% sorts and compresses numeric citations
\usepackage[colorlinks=true,linkcolor=Blue,citecolor=Blue,urlcolor=Blue]{hyperref}

% ---------------------------------------------------------------------
% Shared notation
\newcommand{\R}{\mathbb{R}}
\newcommand{\E}{\mathbb{E}}
\renewcommand{\P}{\mathbb{P}}
\newcommand{\one}{\mathbf{1}}
\newcommand{\D}{\mathcal{D}}
\newcommand{\Lgen}{\mathcal{L}}
\DeclareMathOperator{\Var}{Var}
\DeclareMathOperator{\Cov}{Cov}
\DeclareMathOperator{\Ent}{Ent}
\DeclareMathOperator{\KL}{KL}
\DeclareMathOperator{\Tr}{Tr}
\DeclareMathOperator{\ran}{ran}
\newcommand{\norm}[1]{\lVert #1\rVert}
\newcommand{\ip}[2]{\langle #1,#2\rangle}
% Locators.  \ploc{compiled locator}{TeX file:line} cites the preprint (or a
% companion).  Only the compiled locator (section, equation, result, page of the
% PDF) is printed.  The TeX file:line stays in the source to help checkers find
% the passage and never appears in the PDF.  \sloc{key}{locator} cites a source.
\newcommand{\ploc}[2]{#1}
\newcommand{\sloc}[2]{\cite[#2]{#1}}
% Relationship types and flags, as used in the summary table and the Summary
% paragraphs: \rel{same technique}; \attrgap and \overstated tag a Cited?
% cell (precedence_stats.py checks them against the JSON).
\newcommand{\rel}[1]{\textit{#1}}
\newcommand{\attrgap}{\textbf{attribution gap}}
\newcommand{\overstated}{\textbf{overstated novelty}}
% Uniform list and table styling across all findings.
\setlist{itemsep=2pt,topsep=3pt,parsep=0pt,leftmargin=1.6em}
\setlist[enumerate]{label=\arabic*.}
\setlist[itemize]{label=\textbullet}
\renewcommand{\arraystretch}{1.15}
% Findings are numbered by LaTeX, so nobody types an ID.  In the main file,
% \findinggroup{F} after the Findings \section sets the ID prefix and restarts
% the count; IDs follow the order of the \input lines (by magnitude).  Each
% fragment opens with
%   \finding{<title>}{<slug>}   or   \finding[<short title>]{<title>}{<slug>}
% which prints an unnumbered subsection "F3. <title>", adds "F3. <title>" (or
% the short title) to the table of contents and PDF bookmarks, and does
% \label{f:<slug>}, so \ref{f:<slug>} prints the ID "F3" as a link.  The
% heading carries no LaTeX section number, so a finding never shows two
% numbers (no "3.2 F2").  Refer to another finding as finding~\ref{f:<slug>}.
% Before any \findinggroup the IDs print as ?1, ?2, ... (as in a standalone
% test of one fragment).
\makeatletter
\newcounter{finding}
\newcommand{\finding@prefix}{?}
\renewcommand{\thefinding}{\finding@prefix\arabic{finding}}
\newcommand{\findinggroup}[1]{\gdef\finding@prefix{#1}\setcounter{finding}{0}}
\newcommand{\finding}[3][]{%
  \stepcounter{finding}%
  \protected@edef\finding@id{\thefinding}% fix the ID now: "F3", not a counter reference
  \edef\finding@next{\noexpand\finding@sect{\finding@id}}%
  \finding@next{#1}{#2}{#3}}
% #1 ID (already expanded), #2 short title (may be empty), #3 title, #4 slug.
\newcommand{\finding@sect}[4]{%
  \subsection*{#1. #3}% hyperref anchors \subsection* itself, just above the heading
  \if\relax\detokenize{#2}\relax
    \addcontentsline{toc}{subsection}{#1. #3}%
  \else
    \addcontentsline{toc}{subsection}{#1. #2}%
  \fi
  \def\@currentlabel{#1}% \ref prints the ID; hyperref links it to this heading
  \label{f:#4}}
\makeatother
% ---------------------------------------------------------------------
\space%

\title{Missing citations in\\ ``A finite-entropy separation of microstates and nonmicrostates free entropy''}
\date{October 12, 2026}

\begin{document}
\maketitle

\section*{Summary of findings}

Each row is a citation the preprint is missing: earlier work that already
contains part of the preprint, which the preprint does not cite where that
part appears. ``Cited?'' names the earlier work and says what, if anything, the
preprint cites it for. The exception is a row flagged \emph{overstated
novelty}: there the preprint cites the work, but its abstract or framing
presents the content as new. Terms are defined in Section~\ref{sec:about};
centrality is a best-effort judgment. Locations use the preprint's own
numbering and printed pages.

% Each row starts "\ref{f:<finding>} & \rel{<relationship>} & <centrality> &",
% in that order: \ref prints the finding's ID (F1, F2, ...) as a link, and
% precedence_stats.py reads those three cells. Put \attrgap and \overstated in
% the Cited? cell of a row exactly when the JSON sets the flag. Never type an
% ID: LaTeX assigns them from the main file's order (by magnitude).
% The Technique or result cell, like a title, names the approach or result,
% with no author names or citations.
\begin{center}\footnotesize
\begin{xltabular}{\linewidth}{@{}l >{\raggedright\arraybackslash}p{0.11\linewidth} >{\raggedright\arraybackslash}p{0.10\linewidth} >{\raggedright\arraybackslash}X >{\raggedright\arraybackslash}p{0.17\linewidth} >{\raggedright\arraybackslash}p{0.24\linewidth}@{}}
\toprule
ID & Relationship & Centrality & Technique or result & In the preprint & Cited?\\
\midrule
\endhead
\ref{f:noised-models-finite-entropy} & \rel{same result} & central &
Finite microstates entropy of a tuple with matrix models (here embeddable in $L(\mathbb F_2)$) plus small free semicircular noise &
Lemma~4.2 (p.~9); models (p.~8) &
Voiculescu 1994 \cite{Voiculescu1994-AnaloguesEntropyFisherII}: cited for other content (definition of $\chi$); Voiculescu 1990 \cite{NL-Voiculescu1990}, Jung 2002 \cite{Jung2002-FreeEntropyDimensionLemma}, Jekel 2025 \cite{Jekel2025CMP-InformationGeometryTypesLargeN}: not cited; \attrgap\\
\addlinespace
\ref{f:elementary-heat-flow-transfer} & \rel{same result} & central &
Microstates entropy bounded by the entropy of GUE-evolved microstate models minus the integrated free Fisher information, with the terminal entropies kept &
Lemma~5.1, (13) (p.~10), proof (pp.~10--11); overview (p.~2) &
Jekel--Pi 2023 \cite{JekelPi2023-ElementaryProofChiLeqChiStar}: credited in part (Prop.~6.1, pp.~2, 10; the elementary proof, p.~1); Shlyakhtenko--Tao 2020 \cite{ShlyakhtenkoTao2020-FractionalFreeConvolutionPowers}: credited elsewhere (through Jekel--Pi, p.~1); \overstated\\
\addlinespace
\ref{f:stein-kernel-chaos-lift} & \rel{same technique} & central &
Fisher-deficit bound by lifting higher-order Stein kernels (here free cumulants) into semicircular chaos &
Lemma~3.1, (5) for $r\ge2$ (p.~5); Prop.~3.2 (pp.~6--7) &
Diez 2023 \cite{Diez2023-FreeSteinKernelsHigherOrderDerivatives}, Fathi--Nelson 2016 \cite{FathiNelson2016-FreeSteinKernelsFreeLSI}, Cébron 2018 \cite{Cebron2018-QuantitativeFourthMomentFree}: not cited; \attrgap\\
\addlinespace
\ref{f:commuting-coordinate-entropy-deficit} & \rel{adapted} & central &
Free semicirculars plus a tensor-independent coordinate and small free noise: finite $\chi$, but matrix models that inherit the commutation lose entropy &
\S4 (p.~7); (12) (p.~9); (19), Lemma~6.1, proof of Theorem~1.1 (pp.~12--14) &
Jekel 2025 \cite{Jekel2025CMP-InformationGeometryTypesLargeN}: not cited; \attrgap\\
\addlinespace
\ref{f:block-covering-unitary-net} & \rel{adapted} & central &
Unitary-net covering of matrix tuples with a common block structure &
Lemma~6.1, Gaussian bound (p.~12; proof p.~13) &
Voiculescu 1996 \cite{NL-Voiculescu1996} (as expounded by Biane 2001 \cite{Biane2001-EntropieLibreAlgebresOperateurs} and Mingo--Speicher 2017 \cite{MingoSpeicher2017-FreeProbabilityRandomMatrices}), Dykema 1996 \cite{Dykema1996-TwoApplicationsFreeEntropy}: not cited; Ge 1996 \cite{NL-Ge1996}: credited in part (``antecedents'', p.~2); \attrgap\\
\addlinespace
\ref{f:semicircular-wick-polynomials} & \rel{same result} & supporting &
Orthonormal semicircular Wick polynomials (free Wick products) by signed deletion of disjoint adjacent pairs &
Lemma~3.1 (p.~5) and its proof (p.~6); proof of Prop.~3.2 (p.~7); \S1 (p.~2) &
Effros--Popa 2003 \cite{EffrosPopa2003-FeynmanDiagramsWickProducts}, Kemp--Nourdin--Peccati--Speicher 2010 \cite{KempEtAl2010-WignerChaosFourthMoment}: not cited; \attrgap\\
\addlinespace
\ref{f:semicircular-perturbation-fisher-facts} & \rel{same result} & peripheral &
Conjugate variables of a free semicircular perturbation as the projected rescaled noise &
Lemma~3.1 (pp.~5--6); used on pp.~6, 8--9 and~11 &
Voiculescu 1998 \cite{NL-Voiculescu1998}: cited for other content (definitions)\\
\bottomrule
\end{xltabular}
\end{center}

\clearpage
\space%

\tableofcontents

\section{About this document}\label{sec:about}

\paragraph{The preprint.}
The preprint is \emph{A finite-entropy separation of microstates and
nonmicrostates free entropy}, dated September~25, 2026, with author line
``OpenAI'' \cite{OAI-FiniteEntropySeparation}. It is one of the manuscripts in
the \texttt{openai/math} release, produced by an internal OpenAI model. It
constructs a bounded self-adjoint tuple $X$ in a tracial von Neumann algebra
with $-\infty<\chi(X)\le\chi^*(X)-\tfrac12<\infty$, where $\chi$ is
Voiculescu's microstates free entropy and $\chi^*$ his nonmicrostates free
entropy. Section~\ref{sec:overview} summarizes the paper and fixes the notation
used throughout this document. The preprint cites one companion preprint from
the same release, \cite{OAI-FreeGroupFactors}, as
\texttt{FreeGroupFactors}, for a result that it contrasts with its own and does
not use. It belongs to the same release, so it is not treated as earlier work
here.

\paragraph{What this document reports.}
Each finding is a citation the preprint is missing: part of the preprint's
mathematics (a technique, construction, estimate, proof route or result) that
an identifiable earlier paper already contains, where the preprint does not
cite that paper for it where it is used. ``Earlier'' means publicly available
before the preprint's date. The one exception, overstated novelty, is
described below. Not reported: textbook material in its textbook form; weaker versions of
intermediate statements; and, for content the preprint uses, cases where it
cites a later source that covers that content rather than the original. An
earlier result on the same problem is reported even when later work that the
preprint cites implies it. Findings are ordered by centrality, then by relationship in
the order listed below, then by position in the preprint.

\paragraph{Relationships.}
A finding's relationship is the one between the preprint and the first
earlier work listed for it (the principal), the work with the highest relationship in the
order below; when several hold that relationship, however closely each matches,
the one the others credit or build on, otherwise the earliest. Any other works listed relate to the same content more weakly or
in part; the finding says how, including which work originated the idea when
that is not the first.
\begin{description}[leftmargin=1.5em,font=\normalfont\itshape]
\item[same result.] An earlier paper, or several earlier papers together, prove essentially the same statement.
\item[same technique.] The earlier technique, construction or theory is
reused essentially intact, possibly in a new setting or with new ingredients
around it.
\item[adapted.] The earlier technique's structure is followed, with a change
at its core: a different object, or a key step replaced.
\item[proof shortcut.] A published result would replace or substantially
shorten a large part of the argument.
\item[same route.] An earlier paper proposed the proof strategy, even
heuristically, and the preprint carries it out.
\item[fragment.] One identifiable piece of the earlier technique is reused.
\item[weaker result.] An earlier paper proves a weaker, special-case or
partial-range version of one of the preprint's main results, or of a
statement it presents as a significant contribution.
\end{description}

\paragraph{Centrality.}
Each finding also gets a best-effort judgment of how much of the preprint's
main argument depends on it: \emph{main result} (reserved for earlier work that
covers the main result itself, or most of it), \emph{central},
\emph{supporting} or \emph{peripheral}. An earlier weaker result is never
labelled main result. Each finding states the chain of results it feeds, so
the basis for the judgment is visible.

\paragraph{Citations.}
The ``Cited?'' column of the summary names the earlier work and says what the
preprint cites for the content in question: \emph{not cited} (the work is not
in the bibliography, and nothing the preprint cites credits it for this
content), \emph{cited for other content}, \emph{credited elsewhere} (credited
for this content, directly or through a cited location that credits it, but
not where it is used, in a way that an expert reader who doesn't know the earlier work might not connect with it), or \emph{credited in part}. A citation counts only when it is \emph{for that specific content};
citing the same paper nearby for a different result does not count. Content
the preprint credits where it is used, directly or through a cited location
that itself credits the earlier work, is not reported, except when its novelty is overstated (below). For
overstated novelty, the column says where the preprint does cite the work.

\paragraph{Flags.}
\begin{description}[leftmargin=1.5em,font=\normalfont\bfseries]
\item[Attribution gap.] Marks the most serious missing citations: the
relationship is same technique, adapted or same result; the centrality is main
result, central or supporting; and the preprint cites the principal earlier
work for this content nowhere.
\item[Overstated novelty.] The preprint does cite the earlier work for this
content, but its abstract, a statement of its significant contributions, or the framing of
a main result presents that content as new, explicitly or in the way a reader
would reasonably take it. This is the one case in which content the preprint
credits is reported.
\end{description}

\paragraph{How to check a claim.}
Locations in the preprint use its section,
equation and theorem numbers and the page numbers printed in its PDF.
Locations in other sources give the theorem, equation and printed page of the
version listed in the bibliography. Each finding ends with a short checklist
of the passages to put side by side.

\paragraph{Provenance.}
This document was generated by LLM agents (Claude Opus 5.5) running the
precedence-review skill of jnshi/math-lineage (commit
\href{https://github.com/jnshi/math-lineage/tree/e025345/.agents/skills/precedence-review}{\texttt{e025345}}
for the section passes and commit
\href{https://github.com/jnshi/math-lineage/tree/13ccc19/.agents/skills/precedence-review}{\texttt{13ccc19}}
for the later passes), makes no claim of completeness, and has not been
vetted by a human. It was run by Jonathan Shi, who also supplied the leads
behind findings~\ref{f:noised-models-finite-entropy}--\ref{f:semicircular-perturbation-fisher-facts},
harvested from a run of an earlier version (v1) of the skill on this
preprint, including his own lead behind
finding~\ref{f:commuting-coordinate-entropy-deficit}. This run was also used to
test updates to the skill.
\space%
\section{The preprint in brief, and notation}\label{sec:overview}

\paragraph{Main result.}
\ploc{Theorem~1.1 (p.~2)}{01-introduction.tex:42--51}: there exist $n\ge2$
and a bounded self-adjoint tuple $X=(X_1,\dots,X_n)$ in a von Neumann algebra
with faithful normal tracial state $\tau$ such that
\[
 -\infty<\chi(X)\le\chi^*(X)-\tfrac12<\infty,
\]
where $\chi$ is Voiculescu's microstates entropy with an operator-norm cutoff
and a limsup over matrix sizes. The preprint presents this as a negative answer
to the finite-entropy form of the equality question
$\chi(X)>-\infty\Rightarrow\chi(X)=\chi^*(X)$
\ploc{(\S1, p.~1)}{01-introduction.tex:17--26}.

\paragraph{Entropies \ploc{(\S2, p.~3)}{02-conventions.tex:12--90}.}
$\Gamma_R(X;m,d,\varepsilon)$ is the set of self-adjoint $d\times d$ matrix
tuples with coordinate norms at most $R$ whose normalized traces of words of
length at most $m$ are within $\varepsilon$ of those of $X$;
$\chi_R(X)=\inf_{m,\varepsilon}\limsup_{d}[d^{-2}\log\operatorname{Vol}\Gamma_R(X;m,d,\varepsilon)+\frac n2\log d]$
and $\chi(X)=\sup_R\chi_R(X)$. For a tuple $Z$, the conjugate variables $\xi_i$
satisfy $\tau(\xi_iP(Z))=(\tau\otimes\tau)((\partial_iP)(Z))$ for
noncommutative polynomials $P$, with $\partial_i$ the free difference quotient;
$\Phi^*(Z)=\sum_i\norm{\xi_i}_2^2$, and, with $S$ a free standard semicircular
family,
\[
 \chi^*(Z)=\tfrac n2\log(2\pi e)+\tfrac12\int_0^\infty\Big[\frac n{1+s}-\Phi^*(Z+\sqrt s\,S)\Big]\,ds
 \qquad\ploc{(2), p.~3}{02-conventions.tex:73--78}.
\]
On the matrix side, $\operatorname{tr}_d=d^{-1}\Tr$, a standard GUE tuple
$G^{(d)}$ has independent coordinates whose real Hilbert--Schmidt entries have
variance $1/d$, and the normalized classical entropy of a random tuple with
density $f$ is $H_d=-d^{-2}\int f\log f+\frac n2\log d$.

\paragraph{The signal \ploc{(\S4, p.~7)}{04-signal.tex:8--21}.}
Fix $p\ge2$ and $n=p+1$. $B_1,\dots,B_p$ are free variance-one semicirculars,
and $Y$, in a separate tensor factor, is uniform on the $p$ equally spaced
points $y_l=-1+2(l-1)/(p-1)$ of $[-1,1]$, so $Y$ commutes with every $B_i$.
$C=(B_1,\dots,B_p,Y)$, $c_i=\tau(C_i^2)$, $V_a=C+\sqrt a\,S$,
$v_i(u)=u+c_i$, and $L(a)=\frac12\sum_i\log(2\pi e(a+c_i))$ is the entropy of
the semicircular family (or of the independent GUE law) with the same
coordinate variances. $\kappa_m(C)$ is the order-$m$ free-cumulant array of
$C$. The tuple of Theorem~1.1 is $X=V_h$ for a small $h>0$; the second noise
scale is $t=p^{3/4}$, with $p$ fixed large, as in
\ploc{(11), p.~9}{04-signal.tex:110--125}. The spectral projections of
$X_n=Y+\sqrt h\,S_n$ for the intervals between the midpoints of the $y_l$ are
denoted $e_l$; see \ploc{(12), p.~9}{04-signal.tex:127--147}.

\paragraph{The proof routes, as the preprint names them.}
The introduction says that the proof ``compares two noise scales''
\ploc{(\S1, p.~2)}{01-introduction.tex:58--93}.
\begin{itemize}
\item \emph{The analytic estimate} (\S\S3--4): a squared-cumulant bound on the
Fisher deficit $\Phi^*(V_u)-\sum_i v_i(u)^{-1}$
\ploc{(Proposition~3.2, (6), p.~6)}{03-fisher.tex:100--113}, proved through
the ``Wick testing'' identity of
\ploc{Lemma~3.1 (p.~5)}{03-fisher.tex:28--49}, together with the sparse
cumulants of $C$ \ploc{(Lemma~4.1, p.~8)}{04-signal.tex:45--75}, gives
$\chi^*(V_a)>L(a)-\frac12$ for large $a$, which is
\ploc{(9), p.~9}{04-signal.tex:102--104}. Finiteness,
$\chi(X)>-\infty$, is \ploc{Lemma~4.2 (p.~9)}{04-signal.tex:149--170}.
\item \emph{The transfer between the scales}: an entropy-production lemma for
GUE heat flow started from uniform laws $A^{(d)}$ on near-optimal microstate
sets, \ploc{Lemma~5.1 (p.~10)}{05-flow.tex:10--21}:
the inequality \ploc{(13)}{05-flow.tex:16--19},
$\limsup_d H_d(A^{(d)}+\sqrt t\,G^{(d)})\ge\chi_R(X)+\frac12\int_0^t\Phi^*(V_{h+s})\,ds$,
which the preprint describes as using
``the classical-to-free Fisher comparison of Jekel and Pi''.
\item \emph{The matrix estimate} (\S6): the evolved laws $\mu_d$ of
$A^{(d)}+\sqrt t\,G^{(d)}$ retain a block structure, defined by projection
decompositions $P=(P_l)$ of small rank and the compressions
$\D_PD=(\sum_lP_lD_iP_l)_i$, with probability at least $\frac12$, while the
Gaussian reference law $\gamma_d$ with variances $v_i(h+t)$ has it with
probability at most $e^{-4d^2}$
\ploc{(Lemma~6.1, p.~12)}{06-blocks.tex:79--87}. Relative-entropy data
processing turns this into a classical entropy deficit, and the proof of
Theorem~1.1 combines it with Lemma~5.1 and (9)--(10)
\ploc{(pp.~13--14)}{06-blocks.tex:154--209}.
\end{itemize}
\space%

% IDs are assigned by LaTeX in the order of the \input lines (F1, F2, ...).
% Order by magnitude: centrality (main result, central, supporting,
% peripheral), then the relationship (same result, same technique, adapted,
% proof shortcut, same route, fragment, weaker result), then reading order.
% precedence_stats.py checks this order.
% Withdrawing or reordering a finding is just deleting or moving its \input line.
\section{Findings}\label{sec:findings}
\findinggroup{F}
\finding[Finite entropy of a small free semicircular perturbation]{Finite microstates entropy of a tuple perturbed by small free semicircular noise}{noised-models-finite-entropy}

\paragraph{Summary.}
Voiculescu proved that $\chi(X_1+Y_1,\dots,X_m+Y_m)\ge\chi(Y_1,\dots,Y_m)$
when $Y_1,Y_2,\dots$ are free and of finite entropy and the $X_j$ lie in
$W^*(Y_{m+1},Y_{m+2},\dots)$
\sloc{Voiculescu1994-AnaloguesEntropyFisherII}{Prop.~5.5, p.~431}, and applied
it to a tuple in a free group factor plus free semicircular noise (proof of Prop.~6.5,
p.~435); the preprint cites this paper only for the definition of $\chi$ in
\ploc{\S1 (p.~1)}{01-introduction.tex:5--8}.
\ploc{Lemma~4.2 (p.~9)}{04-signal.tex:149--170}, the finiteness half of
Theorem~1.1, proves $\chi(X)\ge\frac n2\log(2\pi h)$ for $X=C+\sqrt h\,S$,
citing only the preprint's own convergence facts.
\rel{Same result}: as $Y$ is atomic,
$W^*(C)\cong L(\mathbb F_p)\mathbin{\bar\otimes}\ell^\infty_p$ embeds
trace-preservingly in $L(\mathbb F_2)$ by Voiculescu's compression theorem
\cite{NL-Voiculescu1990}, not cited, and Prop.~5.5 then gives
$\chi(X)\ge\frac n2\log(2\pi eh)$, larger by $n/2$. Jung's inequality
$\chi(a+\varepsilon Z:Z)\ge\mathbb P_{2c\varepsilon\sqrt n}(a)+n\log\varepsilon+\underline\chi(Z)$,
for $Z$ free from $a$ and $c=\max_i\norm{Z_i}$
\sloc{Jung2002-FreeEntropyDimensionLemma}{proof of Lemma~2.2, p.~4}, not
cited, gives this bound directly for any $a$ with matrix models, whose packing
term $\mathbb P$ is then nonnegative (\rel{same result}, later than
Voiculescu's). Its proof is the lemma's argument, with noise uniform on
microstate sets where the preprint takes GUE noise: microstates of the
unperturbed tuple plus noise asymptotically free from them lie in the sum's
microstate sets, so the noise alone bounds their volume and the tuple's own
entropy, $-\infty$ here, plays no role; cf.\ Jekel's proof of
\sloc{Jekel2025CMP-InformationGeometryTypesLargeN}{Claim~4.15, pp.~46--47},
not cited, with the same structure for free semicirculars and a
tensor-independent coordinate.

\paragraph{Correspondence.}
\begin{center}\small
\begin{xltabular}{\linewidth}{@{}>{\raggedright\arraybackslash}p{0.07\linewidth}>{\raggedright\arraybackslash}X>{\raggedright\arraybackslash}X>{\raggedright\arraybackslash}X@{}}
\toprule
 & Preprint & Voiculescu \cite{Voiculescu1994-AnaloguesEntropyFisherII} & Jung \cite{Jung2002-FreeEntropyDimensionLemma}\\
\midrule
\endfirsthead
\toprule
 & Preprint & Voiculescu \cite{Voiculescu1994-AnaloguesEntropyFisherII} & Jung \cite{Jung2002-FreeEntropyDimensionLemma}\\
\midrule
\endhead
Tuple & $C$; models in dimensions $pq_k$ \ploc{(p.~8)}{04-signal.tex:23--31} & $X_j\in W^*(Y_{m+1},\dots)$; here $C$, via $x\otimes\delta_l\mapsto E_{ll}\otimes x$ into $M_p(\mathbb C)\otimes L(\mathbb F_{1+p^2})\cong L(\mathbb F_2)$ \sloc{MingoSpeicher2017-FreeProbabilityRandomMatrices}{(6.3), p.~172} & $a$ with matrix models, so $\mathbb P_{2c\varepsilon\sqrt n}(a)\ge0$\\
\addlinespace
Noise & $\sqrt h\,S$, free from $C$, $0<h<1$ & $Y_j=\varepsilon S_j$, free from the free group factor; $\varepsilon=\sqrt h$ & $\varepsilon Z$, $Z$ free from $a$, $\varepsilon<1/(2c+2)$; here $Z=S$, $c=2$, $\varepsilon=\sqrt h<\frac16$\\
\addlinespace
Step & noised models in $\Gamma_{R_0}(X;m,d,\varepsilon)$ with probability $q_d\to1$, density at most $(d/2\pi h)^{nd^2/2}$ \ploc{(p.~9)}{04-signal.tex:153--165} & Props.~2.3, 3.8, 5.4 (pp.~414--430); no matrices & packing centres of $\Gamma_R(a)$ translated by noise microstates $(m,\gamma)$-free from them: disjoint sets in $\Gamma(a+\varepsilon Z:\varepsilon Z)$ (pp.~3--4)\\
\addlinespace
Bound & $\chi(X)\ge\frac n2\log(2\pi h)$ \ploc{(p.~9)}{04-signal.tex:166--169} & $\chi(X)\ge\chi(\sqrt h\,S)=\frac n2\log(2\pi eh)$ & $\chi(X)\ge\chi(X:S)\ge\frac n2\log h+\underline\chi(S)=\frac n2\log(2\pi eh)$\\
\addlinespace
\bottomrule
\end{xltabular}
\end{center}

\paragraph{How to check.}
\begin{enumerate}
\item \ploc{Lemma~4.2 (p.~9)}{04-signal.tex:149--170} against
\sloc{Voiculescu1994-AnaloguesEntropyFisherII}{Prop.~5.5, p.~431}, whose
printed conclusion has $n$ where $m$ is meant, and the display in the proof of
Prop.~6.5 (p.~435) with $\varepsilon=\sqrt h$; the tensor description of $C$
\ploc{(\S4, p.~7)}{04-signal.tex:8--18} against
\sloc{MingoSpeicher2017-FreeProbabilityRandomMatrices}{(6.3), p.~172}.
\item The proof of \ploc{Lemma~4.2}{04-signal.tex:153--169}, with the models of
\ploc{p.~8}{04-signal.tex:23--31}, against the chain of inequalities in Jung's
proof of \sloc{Jung2002-FreeEntropyDimensionLemma}{Lemma~2.2, pp.~3--4} and
Jekel's proof of
\sloc{Jekel2025CMP-InformationGeometryTypesLargeN}{Claim~4.15, p.~47}.
\end{enumerate}

\space%
\finding[Heat-flow transfer from evolved matrix models to microstates entropy]{Microstates entropy bounded by the entropy of GUE-evolved microstate models minus the integrated free Fisher information}{elementary-heat-flow-transfer}

\paragraph{Summary.}
Jekel and Pi's proof of $\chi\le\chi^*$
\sloc{JekelPi2023-ElementaryProofChiLeqChiStar}{proof of Thm~6.4} obtains,
at the foot of p.~27 and for $\mathcal B=\mathbb C$,
the inequality \ploc{(13)}{05-flow.tex:16--19} of
\ploc{Lemma~5.1 (p.~10)}{05-flow.tex:10--21}, keeping the entropies of the
GUE-evolved models and bounding them by the Gaussian maximum only on p.~28. The preprint cites Jekel--Pi for an elementary proof of
$\chi\le\chi^*$ \ploc{(p.~1)}{01-introduction.tex:30--32} and, at the
transfer, only for their Proposition~6.1, the classical-to-free Fisher
comparison \ploc{(pp.~2, 10)}{01-introduction.tex:95--96; 05-flow.tex:3--4};
its overview of the proof of Theorem~1.1
\ploc{(p.~2)}{01-introduction.tex:95--98} presents the finite-time inequality
that retains the terminal entropies as what Lemma~5.1 adds to that comparison.
\rel{same result}: Jekel--Pi prove (13) for the same ensembles, uniform laws
on cutoff microstate sets, along an ultrafilter where the preprint selects one
dimension per accuracy level, and the proof of Lemma~5.1
\ploc{(pp.~10--11)}{05-flow.tex:23--130} follows theirs step by step, proving
two inputs itself \ploc{(Lemmas~2.1--2.2, pp.~3--5)}{02-conventions.tex:36--177}.
Their approach compares $\chi$ with the integrated free Fisher information
without large deviations: realize $\chi_R$ by uniform laws on near-optimal
microstate sets (the realization of $\chi$ by random matrix models formulated
in Jekel's appendix to Shlyakhtenko--Tao
\sloc{ShlyakhtenkoTao2020-FractionalFreeConvolutionPowers}{Prop.~B.7}, which
Jekel--Pi credit), add GUE noise, integrate the classical Fisher information by
de Bruijn's identity, and bound it below in the limit by $\Phi^*$, since matrix
integration by parts turns the classical score tested against polynomials into
the conjugate-variable functional. The preprint's own step comes next: its
\S6 bounds the retained entropies by a relative-entropy deficit
\ploc{(p.~2)}{01-introduction.tex:99--101}.

\paragraph{Correspondence.}
\begin{center}\small
\begin{xltabular}{\linewidth}{@{}>{\raggedright\arraybackslash}p{0.11\linewidth}>{\raggedright\arraybackslash}X>{\raggedright\arraybackslash}X@{}}
\toprule
Step & Preprint & Jekel--Pi \cite{JekelPi2023-ElementaryProofChiLeqChiStar}, $\mathcal B=\mathbb C$\\
\midrule
\endfirsthead
\toprule
Step & Preprint & Jekel--Pi \cite{JekelPi2023-ElementaryProofChiLeqChiStar}, $\mathcal B=\mathbb C$\\
\midrule
\endhead
Models & Uniform on $\Gamma_R(X;k,d_k,1/k)$, log-volume $\ge\chi_R(X)-1/k$, one $d_k$ per $k$; \ploc{(14), p.~10}{05-flow.tex:28--39} & Uniform on $\Gamma_{\mathbf R}^{(n)}(\mathcal O_k)$, $h^{(n)}\ge\chi^{\mathcal U}_{\mathbf R}-1/k$ on $A_k$; proof of Thm~4.10, pp.~18--19; Rem.~4.14, p.~19\\
\addlinespace
Entropy gain & \ploc{(15)--(16), p.~10}{05-flow.tex:41--70}: $\frac{d}{ds}H_d(A^{(d)}+\sqrt s\,G^{(d)})=\frac12I_d(s)$, $I_d$ the normalized Fisher information & (3.2), p.~11: $h^{(n)}(\mathbf X+t^{1/2}\mathbf S)-h^{(n)}(\mathbf X)=\frac12\int_0^t\mathcal I^{(n)}(\mathbf X+u^{1/2}\mathbf S)\,du$\\
\addlinespace
Fisher bound & $\liminf_dI_d(s)\ge\Phi^*(V_{h+s})$ via matrix integration by parts and Lemma~2.1; \ploc{(17)--(18), p.~11}{05-flow.tex:72--115} & Prop.~6.1, p.~24: $\Phi^*(\mathbf X+u^{1/2}\mathbf S)\le\lim_{n\to\mathcal U}\mathcal I^{(n)}(\mathbf X^{(n)}+u^{1/2}\mathbf S^{(n)})$; via Cor.~5.5 and Lemma~5.8, pp.~22--23\\
\addlinespace
Result & \ploc{(13)}{05-flow.tex:16--19}: $\limsup_dH_d(A^{(d)}+\sqrt t\,G^{(d)})\ge\chi_R(X)+\frac12\int_0^t\Phi^*(V_{h+s})\,ds$ & p.~27: $\lim_{n\to\mathcal U}h^{(n)}(\mathbf X^{(n)}+t^{1/2}\mathbf S^{(n)})\ge\chi^{\mathcal U}(\mathbf X)+\frac12\int_0^t\Phi^*(\mathbf X+u^{1/2}\mathbf S)\,du$\\
\addlinespace
\bottomrule
\end{xltabular}
\end{center}

\paragraph{How to check.}
\begin{enumerate}
\item \ploc{(13), p.~10}{05-flow.tex:16--19} against the foot of
\sloc{JekelPi2023-ElementaryProofChiLeqChiStar}{p.~27} with
$\mathcal B=\mathbb C$ and the normalizations of p.~11; then p.~28 and the
overview \ploc{(p.~2)}{01-introduction.tex:95--98}. Lemma~4.9's proof
(pp.~15--16) gives $\mathcal U$ with
$\chi^{\mathcal U}_{\mathbf R}=\chi_{\mathbf R}$ at a cutoff above $\norm{X}$;
indices $n_k$ in the sets $A_k$ of Theorem~4.10's proof (p.~18), approaching
the ultralimit, give (13) as stated.
\item The proof of Lemma~5.1 \ploc{(pp.~10--11)}{05-flow.tex:23--130} step by
step against those of Theorem~6.4 (p.~27) and Proposition~6.1 (pp.~25--26);
Fatou replaces Corollary~3.9 (p.~9) and Lemma~6.3 (p.~26).
\end{enumerate}

\space%
\finding[Higher-order Stein kernels lifted to semicircular chaos]{A Fisher-deficit bound from higher-order Stein kernels lifted to semicircular chaos}{stein-kernel-chaos-lift}

\paragraph{Summary.}
Diez \sloc{Diez2023-FreeSteinKernelsHigherOrderDerivatives}{Lemma~4, p.~29}
bounds the free Fisher information of a semicircular smoothing, relative to a
semicircular law, by the free Stein discrepancy of any order~$k$, carrying to
higher order the lemma of Fathi and Nelson
\sloc{FathiNelson2016-FreeSteinKernelsFreeLSI}{Lemma~2.4, p.~9}, who introduced
free Stein kernels and on whom he builds (p.~29); Cébron
\sloc{Cebron2018-QuantitativeFourthMomentFree}{Lemma~2.6, p.~5} extends their
lemma, for one variable, to signals without regularity.
The preprint cites none of these works, and its citations for the cumulant
calculus and for conjugate variables do not cover this argument.
Its proof of the squared-cumulant bound on the Fisher deficit,
\ploc{Proposition~3.2, (6) (pp.~6--7)}{03-fisher.tex:100--167}, with the lift
\ploc{(5) (p.~5)}{03-fisher.tex:28--44} for $r\ge2$, follows this argument;
through \ploc{(7)--(9) (pp.~8--9)}{04-signal.tex:77--104} it gives
$\chi^*(V_{h+t})>L(h+t)-\frac12$ in the
\ploc{proof of Theorem~1.1 (p.~14)}{06-blocks.tex:194--206}.
The shared approach, the free Stein-kernel argument (Fathi--Nelson; at higher
order, Diez), writes the relative score as the conditional expectation of the
signal's Stein kernel lifted into the semicircular chaos of the noise, an
order-$r$ kernel becoming an order-$r$ Wick polynomial with weight $u^{-r/2}$,
and bounds it by the chaos isometry.
\rel{same technique}: the kernels are the constant cumulant arrays
$\kappa_{r+1}(C)$ of all orders $r\ge2$ at once, and the lift is proved
combinatorially, so, as for Cébron, no regularity of $C$ is needed
($\Phi^*(C)=\infty$, as $Y$ is atomic).

\paragraph{Correspondence.}
Preprint: $Z=V_u$, $\xi_i$ its conjugate variables, $D_{\mathbf j}(P)$ the
trace $\tau^{\otimes(r+1)}$ of the $r$-th iterated free difference quotient of
$P$ at $Z$. Diez: $X$ is the preprint's $C$,
$X(t)=e^{-t}X+\sqrt{1-e^{-2t}}\,\rho^{-1/2}S$ ((100), p.~28), $E_t$ the
conditional expectation onto $W^*(X(t))$, $\Xi(t)$ its conjugate variables,
$A_k$ an order-$k$ Stein kernel; $Z=e^tX(t)$ for $\rho=1$, $u=e^{2t}-1$.
\begin{center}\small
\begin{xltabular}{\linewidth}{@{}>{\raggedright\arraybackslash}p{0.1\linewidth}>{\raggedright\arraybackslash}X>{\raggedright\arraybackslash}X@{}}
\toprule
 & Preprint & Diez \cite{Diez2023-FreeSteinKernelsHigherOrderDerivatives}\\
\midrule
\endfirsthead
\toprule
 & Preprint & Diez \cite{Diez2023-FreeSteinKernelsHigherOrderDerivatives}\\
\midrule
\endhead
Stein identity & $\tau(Z_iP(Z))-v_iD_i(P)=\sum_{r\ge2,\mathbf j}\kappa_{r+1}(C_i,C_{j_1},\dots,C_{j_r})D_{\mathbf j}(P)$ \ploc{(pp.~6--7)}{03-fisher.tex:123--133} & $\ip{\rho X}{P(X)}-(\tau\otimes\tau)\circ\Tr([\mathcal JP](X)^*)=\ip{A_k}{[\mathcal J^kP](X)}_2$ (Def.~17, p.~18)\\
\addlinespace
Lift & $D_{\mathbf j}(P)=u^{-r/2}\tau(U_{\mathbf j}(S)^*P(Z))$ \ploc{(5), p.~5}{03-fisher.tex:42--44} & order-$k$ conjugate variables of $X(t)$: $\rho^{k/2}(1-e^{-2t})^{-k/2}E_t(U_k^{\mathbf j}(S))$ ((103), p.~28)\\
\addlinespace
Residual & $Z_i-v_i\xi_i=E[R_i\mid W^*(Z)]$, $R_i=\sum_{r\ge2,\mathbf j}u^{-r/2}\kappa_{r+1}(C_i,C_{j_1},\dots,C_{j_r})U_{\mathbf j}(S)^*$ \ploc{(p.~7)}{03-fisher.tex:135--157} & $\Phi^*(X(t)\mid V_\rho)=c_t\ip{\partial_S^{k,*}(-A_k)}{\Xi(t)-\rho X(t)}_2$, $c_t=e^{-(k+1)t}\rho^{k/2}(1-e^{-2t})^{-k/2}$ ((123), p.~34)\\
\addlinespace
Norm & $\sum_i\norm{R_i}_2^2=\sum_{r\ge2}u^{-r}\norm{\kappa_{r+1}(C)}_2^2$; contraction and $v_i\ge u$ give (6) \ploc{(p.~7)}{03-fisher.tex:140--166} & $\norm{\partial_S^{k,*}(A_k)}_2=\norm{A_k}_2$ ((118), p.~31; (122), p.~34); Cauchy--Schwarz gives (106)\\
\addlinespace
\bottomrule
\end{xltabular}
\end{center}

\paragraph{How to check.}
\begin{enumerate}
\item $Z_i-v_i\xi_i=E[R_i\mid W^*(Z)]$ \ploc{(p.~7)}{03-fisher.tex:135--157}
against \sloc{Diez2023-FreeSteinKernelsHigherOrderDerivatives}{(107)--(110),
pp.~29--30, and (123), p.~34}, and against Cébron's first-order
$\sqrt{1-e^{-2t}}\,\tau(F\mid F_t)=e^{-t}\tau(\partial_S^*A\mid F_t)$
\sloc{Cebron2018-QuantitativeFourthMomentFree}{p.~6}.
\item \ploc{(5)}{03-fisher.tex:42--44} against
\sloc{Diez2023-FreeSteinKernelsHigherOrderDerivatives}{(103), p.~28} under
$Z=e^tX(t)$, which scales order-$k$ conjugate variables by $e^{-kt}$; at $C=0$,
$u=1$, \ploc{(5)}{03-fisher.tex:42--44} is Diez's definition
$U_k^{\mathbf j}(S)=\partial_S^{k,*}(1^{\otimes(k+1)})$ (Def.~15, p.~17).
\item The \ploc{proof of (6), p.~7}{03-fisher.tex:159--166}, before $v_i\ge u$, against
\sloc{Diez2023-FreeSteinKernelsHigherOrderDerivatives}{(106), p.~29}, applied
formally ($C$ meets neither Lemma~4's hypotheses nor $c_n=1$): for $c_i=1$
both bound $\Phi^*(V_u)-n/(1+u)=(1+u)^{-1}\Phi^*(X(t)\mid V_1)$, by
$(1+u)^{-2}$ times $u^{-k}\Sigma_k^*(C\mid V_1)^2$ and
$\sum_{r\ge2}u^{-r}\norm{\kappa_{r+1}(C)}_2^2$.
\end{enumerate}

\space%
\finding[Commuting coordinate kept through free noise forces an entropy deficit]{A coordinate commuting with free semicirculars, kept through small free noise, forces an entropy deficit on matrix models}{commuting-coordinate-entropy-deficit}

\paragraph{Summary.}
Jekel's counterexample in the setting of laws
\sloc{Jekel2025CMP-InformationGeometryTypesLargeN}{\S4.4, pp.~45--51} perturbs
two free semicirculars and a semicircular tensor-independent of them by small
free semicircular noise, obtaining a law $\mu$ with finite
$\chi^{\mathcal U}$; bounded matrix models of a second, free law $\nu$ that
are $L^2$-close to models of $\mu$ inherit the near-commutation and fall short
of $\chi^{\mathcal U}(\nu)$ by an amount bounded away from zero.
The preprint builds the tuple of its
\ploc{Theorem~1.1}{01-introduction.tex:42--51} on the same construction, in
\ploc{\S4 (p.~7)}{04-signal.tex:3--21}, and proves the same kind of deficit
in \ploc{Lemma~6.1 and the proof of Theorem~1.1 (pp.~12--14)}{06-blocks.tex:79--193},
but does not cite Jekel's paper anywhere; its one nearby citation bearing on
this design is to Ge's argument for prime factors \cite{NL-Ge1996}, given on
\ploc{p.~2}{01-introduction.tex:83--85} only as an antecedent for commutation
constraints on microstate volumes in general.
The shared approach follows Jekel's construction: small free noise makes
$\chi$ finite, while the tensor coordinate's near-commutation survives and
costs entropy to the matrix models that inherit it.
\rel{adapted}: the preprint makes the coordinate atomic, so that
near-commutation becomes block structure, and puts the loss on the law's own
evolved microstate ensembles $\mu_d$, via a block event, a unitary-net
covering and data processing rather than Jekel's conditioning and change of
variables.

\paragraph{Correspondence.}
\begin{center}\small
\begin{xltabular}{\linewidth}{@{}>{\raggedright\arraybackslash}p{0.13\linewidth}>{\raggedright\arraybackslash}X>{\raggedright\arraybackslash}X@{}}
\toprule
 & Preprint & Jekel \cite{Jekel2025CMP-InformationGeometryTypesLargeN}\\
\midrule
\endfirsthead
\toprule
 & Preprint & Jekel \cite{Jekel2025CMP-InformationGeometryTypesLargeN}\\
\midrule
\endhead
Construction & $X=C+\sqrt h\,S$, $h$ small, $Y$ atomic in a separate tensor factor \ploc{(\S4, p.~7)}{04-signal.tex:8--18}; $\chi(X)>-\infty$ \ploc{(Lemma~4.2, p.~9)}{04-signal.tex:149--170} & $\mathbf X=(1-\varepsilon)^{1/2}(S_1,S_2,S_3)+\varepsilon^{1/2}(S_1',S_2',S_3')$: $S_1,S_2$ free, $S_3$ tensor independent of them (pp.~45--46); $\chi^{\mathcal U}(\mu)>-\infty$ (Claim~4.15, pp.~46--47)\\
\addlinespace
Kept after noise & \ploc{(12), p.~9}{04-signal.tex:127--147}: spectral projections $e_l$ of $X_n$ near the $e_l(0)$, which commute with all $B_i$: $\tau(e_l)<3/(2p)$, $\sum_{i,l}\norm{e_lX_ie_l}_2^2>3p/4$ & $\norm{[X_j,X_3]}\le24\varepsilon^{1/2}$, $j=1,2$ (p.~46)\\
\addlinespace
In matrix models & \ploc{(19), p.~12}{06-blocks.tex:12--57}: the midpoint spectral partition $P$ of $A^{(d)}_n$ has $\operatorname{rank}P_l\le2d/p$, $\sum_{i\le p,\,l}\norm{P_lA^{(d)}_iP_l}_{2,d}^2>p/2$ & $\lim_{n\to\mathcal U}\norm{[X_1^{(n)},X_3^{(n)}]}_{L^2}\le24\varepsilon^{1/2}$, passed to $Y_j^{(n)}$ at cost $2R\norm{Y_j^{(n)}-X_j^{(n)}}_{L^2}$ (p.~48)\\
\addlinespace
Entropy deficit & $\limsup_dH_d(A^{(d)}+\sqrt t\,G^{(d)})\le L(h+t)-1$ \ploc{(p.~14)}{06-blocks.tex:189--190}, via a block event, $\mu_d$-probability $\ge1/2$, $\gamma_d$-probability $\le e^{-4d^2}$ \ploc{(Lemma~6.1, p.~12)}{06-blocks.tex:79--87} & for $R$-bounded models $\mathbf Y^{(n)}$ of $\nu$ and $\mathbf X^{(n)}$ of $\mu$: $e^{a/2}[(25+6R)\varepsilon^{1/2}+2Rb^{1/2}]\ge\varepsilon^{1/4}$, $a=\chi^{\mathcal U}(\nu)-\lim_{n\to\mathcal U}h^{(n)}(\mathbf Y^{(n)})$, $b=\lim_{n\to\mathcal U}\norm{\mathbf X^{(n)}-\mathbf Y^{(n)}}_{L^2}^2-d_{W,\mathrm{CEP}}(\mu,\nu)^2$ (Claim~4.16, (4.5), p.~47)\\
\addlinespace
\bottomrule
\end{xltabular}
\end{center}

\paragraph{How to check.}
\begin{enumerate}
\item \ploc{\S4, p.~7}{04-signal.tex:8--21} and
\ploc{(12), p.~9}{04-signal.tex:127--147} against Jekel's construction and
commutator bound \sloc{Jekel2025CMP-InformationGeometryTypesLargeN}{pp.~45--46}.
\item \ploc{(19), p.~12}{06-blocks.tex:12--57} against the passage of the
commutator bound to $\mathbf X^{(n)}$ and $\mathbf Y^{(n)}$
\sloc{Jekel2025CMP-InformationGeometryTypesLargeN}{p.~48}.
\item \ploc{Lemma~6.1 and p.~14}{06-blocks.tex:79--190} against Claim~4.16
\sloc{Jekel2025CMP-InformationGeometryTypesLargeN}{pp.~47--51}: the
preprint's deficit falls on the law's own ensembles, Jekel's on models of
$\nu$ (p.~45), by different estimates.
\end{enumerate}

\space%
\finding[Unitary-net covering of block-structured tuples]{Bounding the measure of matrix tuples with a common block structure by a unitary net over fixed blocks}{block-covering-unitary-net}

\paragraph{Summary.}
Voiculescu proved in the third paper of his free-entropy series
\cite{NL-Voiculescu1996} (FE~III below) that finite free entropy excludes
property~$\Gamma$ and Cartan subalgebras, by a covering argument expounded by
Biane \sloc{Biane2001-EntropieLibreAlgebresOperateurs}{Lemme~6.3,
pp.~289--290} and by Mingo and Speicher
\sloc{MingoSpeicher2017-FreeProbabilityRandomMatrices}{\S7.7.1, Lemma~9,
pp.~196--198}; Dykema's many-block version
\sloc{Dykema1996-TwoApplicationsFreeEntropy}{Lemma~2.2, pp.~3--7} is the
closest in form.
The preprint cites none of these works: its \S6 cites nothing, and its
introduction \ploc{(\S1, p.~2)}{01-introduction.tex:83--87} says only that
controlling microstate volumes by commutation and projection geometry has
``antecedents'' in Ge's argument for prime factors \cite{NL-Ge1996}, and cites
Szarek \cite{Szarek1997-MetricEntropyHomogeneousSpaces} for the net bound.
The Gaussian half of \ploc{Lemma~6.1 (p.~12)}{06-blocks.tex:79--87}, proved
on \ploc{p.~13}{06-blocks.tex:112--151}, adapts FE~III's approach to show that
under $\gamma_d$ the first $p$ coordinates rarely put a fixed fraction of
their energy into the blocks of some decomposition into $p$ projections of
rank at most $2d/p$: conjugate the decomposition to fixed blocks, replace the
unitary by a point of an operator-norm net paid for once, and bound the
measure at each net point, where the gain accumulates over the coordinates.
The insight, that one common unitary makes the constraint strong (Mingo and
Speicher, p.~198), is Voiculescu's; his theorem
\sloc{MingoSpeicher2017-FreeProbabilityRandomMatrices}{Theorem~7(i), p.~195}
already gives $\chi(C)=-\infty$, as $Y$ is central in $W^*(C)$.
\rel{adapted}: the reduction is followed step by step, but the estimate at
each net point becomes a chi-square tail under $\gamma_d$, not a volume of
balls with a thin off-diagonal factor as $\omega\to0$, and the event ranges
over all rank lists.

\paragraph{Correspondence.}
Preprint: $D_{[p]}=(D_1,\dots,D_p)$,
$\norm{D}_{2,d}^2=\sum_i\operatorname{tr}_d(D_i^*D_i)$, $v=v_i(h+t)$ for
$i\le p$, and $P^0$ the standard coordinate blocks of a rank list. Biane: $d$
variables, a microstate $(M_1,\dots,M_d)$ of size $N$, and a projection $\Pi$
of rank $q=[N\tau(P)]$ with $\norm{[M_i,\Pi]}_2<2\omega$.
\begin{center}\small
\begin{xltabular}{\linewidth}{@{}>{\raggedright\arraybackslash}p{0.10\linewidth}>{\raggedright\arraybackslash}X>{\raggedright\arraybackslash}X@{}}
\toprule
 & Preprint, proof of Lemma~6.1 & FE~III, after Biane\\
\midrule
\endfirsthead
\toprule
 & Preprint, proof of Lemma~6.1 & FE~III, after Biane\\
\midrule
\endhead
Fixed blocks & each decomposition is a conjugate $UP^0U^*$; compression
$\D_{P^0}$ \ploc{(p.~13)}{06-blocks.tex:119--120} &
$\Pi=U\operatorname{diag}(I_q,0)U^*$; off-diagonal blocks $C_i$ of
$U^*M_iU$ with $\frac2N\Tr(C_iC_i^*)\le4\omega^2$ (p.~289)\\
\addlinespace
Net & operator-norm $\delta$-net of $U(d)$, $(1+2/\delta)^{2d^2}$ points,
$\delta=1/(64\sqrt v)$ \ploc{(p.~13)}{06-blocks.tex:112--118} &
$(C\delta^{-1})^{N^2-q^2-(N-q)^2}$ points $U_s$ of a Szarek net of
$U(N)/(U(q)\times U(N-q))$ (p.~290)\\
\addlinespace
Conjuga\-tion loss &
$\norm{U^*D_{[p]}U-V^*D_{[p]}V}_{2,d}\le2\delta\norm{D_{[p]}}_{2,d}$, at most
$\sqrt p/8$ on the event \ploc{(p.~13)}{06-blocks.tex:121--133} & commutator
at most $2\omega+8\delta$ at the net point (p.~290)\\
\addlinespace
Per net point &
$\gamma_d(\norm{\D_{P^0}(U^*D_{[p]}U)}_{2,d}^2\ge p/64)\le
e^{d^2(\log2-p/(256v))}$ by chi-square on $p\sum_l\operatorname{rank}(P^0_l)^2\le2d^2$
of $pd^2$ dimensions; union over $(d+1)^p$ rank lists
\ploc{(p.~13)}{06-blocks.tex:134--151} & volume of
$\bigl(U_s^*(B(\sqrt{4N})\times B((2\omega+8\delta)\sqrt N)\times
B(\sqrt{4N}))U_s\bigr)^d$ over $s\in S$, with $\delta=\omega$, giving
$\chi\le C_1+C_2\log\omega$ (p.~290)\\
\addlinespace
\bottomrule
\end{xltabular}
\end{center}

\paragraph{How to check.}
\begin{enumerate}
\item The Gaussian bound's proof \ploc{(p.~13)}{06-blocks.tex:112--151}
against the proof of Lemme~6.3
\sloc{Biane2001-EntropieLibreAlgebresOperateurs}{pp.~289--290} and
\sloc{MingoSpeicher2017-FreeProbabilityRandomMatrices}{pp.~196--198}: in both,
the net is paid for once, against a gain from all $p$, respectively $d\ge2$,
coordinates.
\item The introduction's citations of Ge and Szarek
\ploc{(\S1, p.~2)}{01-introduction.tex:83--87} against
\sloc{MingoSpeicher2017-FreeProbabilityRandomMatrices}{p.~195}, which credits
Theorem~7(i)--(ii) to Voiculescu and (iii), primeness, to Ge.
\end{enumerate}

\space%
\finding{Orthonormal semicircular Wick polynomials given by deleting adjacent pairs}{semicircular-wick-polynomials}

\paragraph{Summary.}
Effros and Popa gave an explicit formula for the $q$-Wick product; at $q=0$ it
is a sum over gap-free diagrams, that is, sets of disjoint adjacent pairs
\sloc{EffrosPopa2003-FeynmanDiagramsWickProducts}{(12), p.~12}. Kemp, Nourdin,
Peccati and Speicher (KNPS), following Biane and Speicher \cite{NL-BianeSpeicher1998},
state that $Y\mapsto Y\Omega$ is an isometry into the full Fock space sending
products of Chebyshev polynomials over runs of equal indices to orthonormal
elementary tensors
\sloc{KempEtAl2010-WignerChaosFourthMoment}{Prop.~3.1 and (3.2), p.~29}. The
preprint cites neither: it names ``orthogonality of semicircular Wick
polynomials'' as a tool \ploc{(\S1, p.~2)}{01-introduction.tex:68--69}, citing
Speicher and Nica--Speicher there only for the cumulant calculus, defines
$U_{\mathbf j}(S)$ by deleting disjoint adjacent pairs
\ploc{(proof of Lemma~3.1, p.~6)}{03-fisher.tex:52--56}, proves its
orthonormality \ploc{(Lemma~3.1, p.~5)}{03-fisher.tex:40--41} by a cumulant
cancellation, and uses it in the
\ploc{proof of Proposition~3.2 (p.~7)}{03-fisher.tex:135--144}.
\rel{same result}: the definition is Effros and Popa's formula at $q=0$ term by
term; orthonormality, which they do not state, follows at once from their
characterization $u\Omega=f_1\otimes\cdots\otimes f_n$ of the Wick product
(p.~7), defined following Bo\.zejko, K\"ummerer and Speicher
\cite{NL-BozejkoKummererSpeicher1997}; and since pairs that straddle two runs
carry $\delta=0$, $U_{\mathbf j}$ is the KNPS product over runs, which on a
single run is the classical expansion of $U_m$.

\paragraph{Correspondence.}
\begin{center}\small
\begin{xltabular}{\linewidth}{@{}>{\raggedright\arraybackslash}p{0.15\linewidth}>{\raggedright\arraybackslash}X>{\raggedright\arraybackslash}X@{}}
\toprule
 & Preprint & Effros--Popa (EP) \cite{EffrosPopa2003-FeynmanDiagramsWickProducts}; KNPS \cite{KempEtAl2010-WignerChaosFourthMoment}\\
\midrule
\endfirsthead
\toprule
 & Preprint & Effros--Popa (EP) \cite{EffrosPopa2003-FeynmanDiagramsWickProducts}; KNPS \cite{KempEtAl2010-WignerChaosFourthMoment}\\
\midrule
\endhead
Family & $U_{\mathbf j}(S)$, $\mathbf j=(j_1,\dots,j_r)$: sum over sets $F$ of disjoint adjacent pairs of the word $S_{j_1}\cdots S_{j_r}$ with the pairs in $F$ deleted, each with factor $-\delta_{j_aj_{a+1}}$ \ploc{(proof of Lemma~3.1, p.~6)}{03-fisher.tex:52--56} & ${:}\xi_1\cdots\xi_r{:}=\sum_{\gamma\in\mathsf{GF}(r)}(-1)^{\#(\gamma)}v(\gamma)$, $v(\gamma)$ the covariances of the pairs times the word of the singletons (EP (12), p.~12; $v$, p.~5); $\xi_a=\varphi(e_{j_a})$ gives $\E(\xi_a\xi_{a+1})=\delta_{j_aj_{a+1}}$\\
\addlinespace
Runs & runs $k_1^{m_1},\dots,k_\ell^{m_\ell}$: $U_{\mathbf j}(S)=U_{m_1}(S_{k_1})\cdots U_{m_\ell}(S_{k_\ell})$ (not stated); e.g.\ $U_{(1,1,2)}=S_1^2S_2-S_2=U_2(S_1)U_1(S_2)$ & $U_{n_1}(X(h_{k_1}))\cdots U_{n_r}(X(h_{k_r}))$, $k_i\ne k_{i+1}$ (their $r$ counts runs; KNPS (3.2), p.~29), $U_n$ as in (2.11), p.~28\\
\addlinespace
Orthonormality & $\tau(U_{\mathbf j}^*U_{\mathbf k})=\delta_{\mathbf j\mathbf k}$ \ploc{(Lemma~3.1, p.~5; proof at $C=0$, $u=1$, p.~6)}{03-fisher.tex:40--41, 90--95} & $u\Omega=f_1\otimes\cdots\otimes f_r$ (EP p.~7), so $\E(u^*v)=\ip{v\Omega}{u\Omega}$; $Y\mapsto Y\Omega$ isometric (KNPS Prop.~3.1, p.~29)\\
\addlinespace
\bottomrule
\end{xltabular}
\end{center}

\paragraph{How to check.}
\begin{enumerate}
\item The first paragraph of the \ploc{proof of Lemma~3.1 (p.~6)}{03-fisher.tex:52--56}
against \sloc{EffrosPopa2003-FeynmanDiagramsWickProducts}{(12) and (13),
p.~12} with $\xi_a=S_{j_a}$: a diagram is gap free when every pair is
$\{a,a+1\}$ (gap $g$, p.~4; $\mathsf{GF}$, p.~11).
\item Orthonormality in \ploc{Lemma~3.1 (p.~5)}{03-fisher.tex:40--41} and
\ploc{its proof (p.~6)}{03-fisher.tex:90--95} against
$u\Omega=f_1\otimes\cdots\otimes f_n$
\sloc{EffrosPopa2003-FeynmanDiagramsWickProducts}{p.~7} at $q=0$ with
orthonormal $f_a$.
\item Factor $U_{\mathbf j}$ over runs and compare with
\sloc{KempEtAl2010-WignerChaosFourthMoment}{Prop.~3.1 and (3.2), p.~29}. On one run of length $m$ the deletion rule gives
$\sum_i(-1)^i\binom{m-i}{i}x^{m-2i}$, the classical expansion of $U_m(x)$.
\end{enumerate}

\space%
\finding[Conjugate variables under free semicircular noise]{Conjugate variables of a free semicircular perturbation as the projected rescaled noise}{semicircular-perturbation-fisher-facts}

\paragraph{Summary.}
Voiculescu's fifth paper on free entropy and Fisher information (FE~V) shows
that, for a free semicircular family $S$ free from $X$, the conjugate
variables of $X+\sqrt\varepsilon\,S$ always exist and equal
$\varepsilon^{-1/2}E[s_i\mid W^*(X+\sqrt\varepsilon\,S)]$
\sloc{NL-Voiculescu1998}{Cor.~3.9}; Major and Shlyakhtenko
\sloc{MajorShlyakhtenko2025-StrongOneBoundedNonMicrostates}{\S2.1, p.~4} and
C\'ebron \sloc{Cebron2018-QuantitativeFourthMomentFree}{p.~6} cite it under
that number, and Fathi and Nelson cite its Ornstein--Uhlenbeck form as FE~V
Prop.~3.7 \sloc{FathiNelson2016-FreeSteinKernelsFreeLSI}{p.~7}. The preprint
cites FE~V only in its introduction, for the definitions of conjugate
variables, free Fisher information and $\chi^*$, and states the formula
without citation as the last statement of
\ploc{Lemma~3.1 (p.~5)}{03-fisher.tex:45--48}, proved as the $r=1$ case of
its Wick-testing identity \ploc{(p.~6)}{03-fisher.tex:95--97}.
\rel{same result}: in the preprint's notation ($X=C$, $\varepsilon=u$,
$Z=C+\sqrt u\,S$) the statement is Voiculescu's formula, whose content is that
the rescaled noise $u^{-1/2}S_i$ already satisfies the conjugate relations of
$Z$, so that its projection onto $L^2(W^*(Z))$ is the conjugate variable for
every $C$; the preprint checks these relations with noncrossing cumulants.
The formula supplies the conjugate variables in the proof of
\ploc{Proposition~3.2 (p.~6)}{03-fisher.tex:117--118}, where it yields
$\Phi^*(V_u)\le n/u$, which makes the integrals in
\ploc{(8)}{04-signal.tex:86--97} and \ploc{(13)}{05-flow.tex:16--19} finite
\ploc{(pp.~8--9 and~11)}{04-signal.tex:90--97, 05-flow.tex:126--127}; its
existence statement lets Lemma~2.1 give \ploc{(18)}{05-flow.tex:110--114} in
the proof of \ploc{Lemma~5.1 (p.~11)}{05-flow.tex:104--114}. Both uses lead to
Theorem~1.1.

\paragraph{Correspondence.}
\begin{center}\small
\begin{xltabular}{\linewidth}{@{}>{\raggedright\arraybackslash}p{0.13\linewidth}>{\raggedright\arraybackslash}X>{\raggedright\arraybackslash}X@{}}
\toprule
 & Preprint & FE~V, as quoted\\
\midrule
\endfirsthead
\toprule
 & Preprint & FE~V, as quoted\\
\midrule
\endhead
Setting & $C$ a centered bounded self-adjoint $n$-tuple, $S$ a standard semicircular family free from $C$, $Z=C+\sqrt u\,S$, $u>0$ \ploc{(Lemma~3.1, p.~5)}{03-fisher.tex:29--31} & $X=(x_1,\dots,x_d)$, $N\subseteq M$ a unital $*$-subalgebra, $S=(s_1,\dots,s_d)$ a free $(0,1)$-semicircular family, free from $N\cup X$; $\varepsilon>0$ \sloc{MajorShlyakhtenko2025-StrongOneBoundedNonMicrostates}{p.~4}\\
\addlinespace
Formula & $\xi_i=u^{-1/2}E[S_i\mid W^*(Z)]$ \ploc{(p.~5)}{03-fisher.tex:45--48} & $\mathcal J(x_i+\sqrt\varepsilon\,s_i:N)=\varepsilon^{-1/2}E_\varepsilon(s_i)$, with $E_\varepsilon$ the conditional expectation onto $W^*(N,X+\sqrt\varepsilon\,S)$ \sloc{MajorShlyakhtenko2025-StrongOneBoundedNonMicrostates}{p.~4}\\
\addlinespace
Existence & the conjugate variables of $Z$ exist, for every such $C$ \ploc{(p.~5)}{03-fisher.tex:45--48} & for any self-adjoint $x_1,\dots,x_n$ and $x_j(t)=e^{-t}x_j+\sqrt{1-e^{-2t}}\,\rho^{-1/2}s_j$, the conjugate variables of $X(t)$ always exist and $\xi_j(t)=\sqrt\rho\,(1-e^{-2t})^{-1/2}\mathcal E_t(s_j)$ \sloc{FathiNelson2016-FreeSteinKernelsFreeLSI}{p.~7}\\
\addlinespace
\bottomrule
\end{xltabular}
\end{center}

\paragraph{How to check.}
\begin{enumerate}
\item The last display of \ploc{Lemma~3.1 (p.~5)}{03-fisher.tex:45--48}
against \sloc{MajorShlyakhtenko2025-StrongOneBoundedNonMicrostates}{\S2.1, p.~4}
with $N=\mathbb C$, $X=C$ and $\varepsilon=u$: then
$\mathcal J(x_i+\sqrt\varepsilon\,s_i:\mathbb C)$ is the $i$-th conjugate
variable $\xi_i$ of the tuple $Z$, and $E_\varepsilon=E[\,\cdot\mid W^*(Z)]$.
\item The same display against
\sloc{FathiNelson2016-FreeSteinKernelsFreeLSI}{p.~7}: $x_j(t)=e^{-t}(x_j+\sqrt\varepsilon\,s_j)$
with $\varepsilon=(e^{2t}-1)/\rho$, and conjugate variables scale as
$\xi_j(\lambda Z)=\lambda^{-1}\xi_j(Z)$; with $\lambda=e^{-t}$ their formula
becomes $\xi_j(X+\sqrt\varepsilon\,S)=\varepsilon^{-1/2}E[s_j\mid W^*(X+\sqrt\varepsilon\,S)]$.
\item The last sentence of the proof of Lemma~3.1
\ploc{(p.~6)}{03-fisher.tex:95--97}: for $r=1$, $U_i(S)=S_i$ and
$D_i(P)=(\tau\otimes\tau)(\partial_iP(Z))$, so \ploc{(5)}{03-fisher.tex:42--44} says
$\tau(u^{-1/2}S_i\,P(Z))=(\tau\otimes\tau)(\partial_iP(Z))$, that is, the
rescaled noise satisfies the conjugate relations of $Z$; projecting it onto
$L^2(W^*(Z))$ gives the right-hand side in item~1.
\end{enumerate}

\space%

\begin{thebibliography}{99}
\small
\bibitem{Voiculescu1994-AnaloguesEntropyFisherII} D. Voiculescu, The analogues of entropy and of Fisher's information measure in free probability theory, II, Invent. Math. 118 (1994) 411--440. doi:10.1007/BF01231539.
\bibitem{NL-Voiculescu1990} D. Voiculescu, Circular and semicircular systems and free product factors, in: Operator Algebras, Unitary Representations, Enveloping Algebras, and Invariant Theory (Paris, 1989), Progr. Math. 92, Birkh\"auser, Boston (1990), 45--60.
\bibitem{Jung2002-FreeEntropyDimensionLemma} K. Jung, A free entropy dimension lemma, Pacific J. Math. 211 (2003), no. 2 (page numbers refer to the arXiv version). arXiv:math/0207149v2.
\bibitem{Jekel2025CMP-InformationGeometryTypesLargeN} D. Jekel, Information geometry for types in the large-n limit of random matrices, Comm. Math. Phys. 406 (2025), no. 11, Art. 272. arXiv:2501.00703v2; doi:10.1007/s00220-025-05451-x.
\bibitem{JekelPi2023-ElementaryProofChiLeqChiStar} D. Jekel, J. Pi, An elementary proof of the inequality chi <= chi* for conditional free entropy, Doc. Math. 29 (2024) 1085--1124 (page numbers refer to the arXiv version). arXiv:2305.02574v4; doi:10.4171/DM/969.
\bibitem{ShlyakhtenkoTao2020-FractionalFreeConvolutionPowers} D. Shlyakhtenko, T. Tao, Fractional free convolution powers, with an appendix by D. Jekel, preprint (2020). arXiv:2009.01882v3.
\bibitem{Diez2023-FreeSteinKernelsHigherOrderDerivatives} C.-P. Diez, Free Stein kernels, free moment maps, and higher order derivatives, preprint 2023. arXiv:2306.00967v1.
\bibitem{FathiNelson2016-FreeSteinKernelsFreeLSI} M. Fathi, B. Nelson, Free Stein kernels and an improvement of the free logarithmic Sobolev inequality, Adv. Math. 317 (2017) 193--223 (page numbers refer to the arXiv version). arXiv:1605.05828v1.
\bibitem{Cebron2018-QuantitativeFourthMomentFree} G. Cébron, A quantitative fourth moment theorem in free probability theory, Adv. Math. 380 (2021) 107579 (page numbers refer to the arXiv version). arXiv:1803.09669v2.
\bibitem{NL-Voiculescu1996} D. Voiculescu, The analogues of entropy and of Fisher's information measure in free probability theory III: The absence of Cartan subalgebras, Geom. Funct. Anal. 6 (1996) 172--199, doi:10.1007/BF02246772
\bibitem{Biane2001-EntropieLibreAlgebresOperateurs} P. Biane, Entropie libre et algèbres d'opérateurs, Séminaire Bourbaki 2000/2001, exp. 889, Astérisque 282 (2002) 279--300.
\bibitem{MingoSpeicher2017-FreeProbabilityRandomMatrices} J. A. Mingo, R. Speicher, Free Probability and Random Matrices, Fields Institute Monographs 35, Springer, 2017.
\bibitem{Dykema1996-TwoApplicationsFreeEntropy} K. Dykema, Two applications of free entropy, Math. Ann. 308 (1997) 547-558 (page numbers refer to the arXiv version). arXiv:funct-an/9605005v2.
\bibitem{NL-Ge1996} L. Ge, Prime factors, Proc. Natl. Acad. Sci. USA 93 (1996), no. 23, 12762--12763, doi:10.1073/pnas.93.23.12762
\bibitem{EffrosPopa2003-FeynmanDiagramsWickProducts} E. G. Effros, M. Popa, Feynman diagrams and Wick products associated with q-Fock space, Proc. Natl. Acad. Sci. USA 100 (2003), no. 15, 8629--8633 (page numbers refer to the arXiv version). arXiv:math/0303045v1; doi:10.1073/pnas.1531460100.
\bibitem{KempEtAl2010-WignerChaosFourthMoment} T. Kemp, I. Nourdin, G. Peccati, R. Speicher, Wigner chaos and the fourth moment, Ann. Probab. 40 (2012), no. 4, 1577--1635 (page numbers refer to the arXiv version). arXiv:1009.3949v2.
\bibitem{NL-Voiculescu1998} D. Voiculescu, The analogues of entropy and of Fisher's information measure in free probability theory. V. Noncommutative Hilbert transforms, Invent. Math. 132 (1998), 189--227, doi:10.1007/s002220050222
\bibitem{OAI-FiniteEntropySeparation} OpenAI, \emph{A finite-entropy separation of microstates and nonmicrostates free entropy}, OpenAI Math Release preprint, September 25, 2026, \url{https://github.com/openai/math}.
\bibitem{OAI-FreeGroupFactors} OpenAI, \emph{An isomorphism of the free group factors}, OpenAI Math Release preprint, September 23, 2026 (cited by the preprint as \texttt{FreeGroupFactors}).
\bibitem{Szarek1997-MetricEntropyHomogeneousSpaces} S. J. Szarek, Metric entropy of homogeneous spaces, Banach Center Publ. 43 (1998) 395--410 (page numbers refer to the arXiv version). arXiv:math/9701213v1.
\bibitem{NL-BianeSpeicher1998} P. Biane, R. Speicher, Stochastic calculus with respect to free Brownian motion and analysis on Wigner space, Probab. Theory Related Fields 112 (1998) 373--409, doi:10.1007/s004400050194
\bibitem{NL-BozejkoKummererSpeicher1997} M. Bo\.zejko, B. K\"ummerer, R. Speicher, q-Gaussian processes: non-commutative and classical aspects, Comm. Math. Phys. 185 (1997) 129--154, doi:10.1007/s002200050084
\bibitem{MajorShlyakhtenko2025-StrongOneBoundedNonMicrostates} B. Major, D. Shlyakhtenko, Establishing strong 1-boundedness via non-microstates free entropy techniques, preprint 2025. arXiv:2510.07558v2.
\end{thebibliography}
\space%

\end{document}
